\documentclass{ximera}
% xmPreamble.tex en xmPrintstyle.sty worden automatisch geladen door ximera.cls (v2.7.8);
% NIET nog eens \input{./preamble.tex} of \addPrintStyle{.} doen (geeft 'already defined'-fouten).

\begin{document}
\author{Alexander Holvoet}
\xmtitle{Een hoop zand op een rechthoekig grondvlak}{}

\begin{exercise}
    Welke zandhoop verwacht je bij een rechthoekig grondvlak? Maak een schets van het te
    verwachten bovenaanzicht en voer daarna pas het experiment uit.
    \begin{oplossing}
        Je krijgt geen piramide meer. Er is niet \'e\'en top maar een lijnstuk met hoogste
        punten. (In de bouw kom je dit soort ruimtelijk lichaam tegen in de vorm van een
        \textit{schilddak}.)
    \end{oplossing}
\end{exercise}

\begin{exercise}
    Wat kun je zeggen over de rusthoek?
    \begin{oplossing}
        Deze is weer constant en gelijk aan die bij de kegel en de piramide met vierkant
        grondvlak. De rusthoek hangt immers enkel af van het soort materiaal.
    \end{oplossing}
\end{exercise}

Hoogtelijnen ken je uit de aardrijkskundeles: het zijn lijnen op een kaart die alle punten
met een gelijke hoogte boven het zeeniveau verbinden.

\begin{exercise}
    Vul je tekening van het bovenaanzicht aan met enkele hoogtelijnen.
    \begin{oplossing}
        De hoogtelijnen zijn evenwijdig met de zijden van de rechthoek en de afstand tussen
        twee opeenvolgende hoogtelijnen is altijd dezelfde. Dit komt omdat de flanken overal
        even steil zijn.
    \end{oplossing}
\end{exercise}

% TODO (afbeelding ontbreekt): in de bron stond hier \afbeelding[6cm]{LRHhoogtelijnen.png}.

De hoogtelijn met de hoogste punten noemen we \textemdash{} net zoals bij echte bergen
\textemdash{} de \emph{kamlijn}. Ze is evenwijdig met de langste zijde van het grondvlak
en staat loodrecht op de kortste zijde. Je ziet dat die in het midden ligt. Het kan zijn
dat je die niet getekend hebt. Verderop proberen we de positie van de kamlijn meetkundig
te bepalen.

\begin{exercise}
    Verbind de hoekpunten van de hoogtelijnen.
    \begin{oplossing}
        Je krijgt een figuur waarin door elk hoekpunt van de hoogtelijnen een schuine
        rechte loopt onder \(45^\circ\) met de zijden.
    \end{oplossing}
\end{exercise}

% TODO (afbeelding ontbreekt): in de bron stond hier \afbeelding[6cm]{LRHhoogtelijnenbiss.png}.

\begin{exercise}
    Aan welke eigenschap voldoen
    \begin{enumerate}[label=(\alph*)]
        \item punten die op de kamlijn liggen, zoals het punt \(R\);
        \item punten op \'e\'en van de schuine rechten die de hoekpunten van de
            rechthoekige hoogtelijnen verbinden, zoals het punt \(P\);
        \item punten die op zowel de kamlijn als op de twee schuine rechten liggen, zoals
            het punt \(S\)?
    \end{enumerate}
    \begin{hint}
        Denk telkens aan de afstand van het punt tot de randen van het grondvlak. Hoe ver
        ligt een zandkorrel op dit punt van de verschillende randen?
    \end{hint}
    \begin{oplossing}
        \begin{enumerate}[label=(\alph*)]
            \item Punten op de kamlijn bevinden zich even ver van de twee overstaande
                zijden, zowel in 3D als in het bovenaanzicht. Dit verklaart waarom de
                kamlijn in het bovenaanzicht in het midden tussen de twee lange randen ligt.
            \item Deze punten liggen even ver van twee opeenvolgende randen, ook weer zowel
                in 3D als in bovenaanzicht. Dit verklaart waarom deze lijnen een hoek van
                \(45^\circ\) maken met deze zijden.
            \item Deze punten bevinden zich even ver van drie opeenvolgende randen.
        \end{enumerate}
    \end{oplossing}
\end{exercise}

% TODO (afbeelding ontbreekt): in de bron stond hier \afbeelding[6cm]{LRHbiss.png}
% (met de punten R, P en S erop aangeduid).

In het bovenaanzicht is de \textit{kamlijn} dus de meetkundige plaats van de punten die
zich even ver bevinden van de twee dichtstbijzijnde zijden. De punten op het uiteinde van
de kamlijn liggen op gelijke afstand van drie zijden. In 3D kunnen we de ribben daarom ook
kamlijnen noemen.

\end{document}
